\begin{figure*}[t]
\centering
\includegraphics[width=\textwidth]{fig_overview_pipeline}
\caption{Overview of the G-theory evaluation pipeline. A factorial design crossing seven facets across eight models and five benchmarks produces 2.3M binary records. G-study variance decomposition via Henderson Method~I reveals that model$\times$item interaction accounts for up to 36\% of variance. D-study projections translate variance components into design-dependent item budgets: a 3-prompt $\times$ 3-seed design achieves $G \geq 0.80$ with 178 items (3.1$\times$ reduction). Three contributions emerge: reproducibility budgets (C1), dominant model$\times$item interaction (C2), and format-associated variance patterns (C3).}
\label{fig:overview}
\end{figure*}

\section{Introduction}
\label{sec:introduction}

LLM benchmark scores determine which models get deployed, which get funded, and which top the public leaderboards---yet their reliability is rarely questioned. It should be. Prompt formatting alone can shift accuracy by up to 76 percentage points \citep{sclar2024formatting, brittlebench2026}; answer ordering \citep{pezeshkpour2024ordering}, numerical precision \citep{yuan2025layercast}, and random seeds \citep{dodge2020finetuning, bui2025seeds} each introduce further instability. These sources have been studied in isolation, but isolation masks the real question: when all sources vary simultaneously, which ones dominate---and does the answer change with the benchmark?

Existing approaches offer partial answers that do not compose. \citet{bouthillier2021accounting} established ANOVA-based variance decomposition for ML reproducibility but did not address LLM evaluation. \citet{messing2026hidden} brought G-theory to LLM evaluation, decomposing five methodology facets---but on a single benchmark, without infrastructure facets or cross-format comparison. IRT-based approaches \citep{polo2024tinybenchmarks, lalor2019irt, zhou2026irt} optimize which items to include but treat prompt, seed, and ordering as fixed conditions, leaving their variance unmodeled. The result is a fragmented picture: we know each source matters individually, but have no unified account of their relative magnitudes, their interactions, or how they shift across benchmark formats.

We provide that unified account. Applying G-theory \citep{shavelson1991generalizability, brennan2001generalizability} to a factorial design crossing seven facets across eight open-weight 7--9B models and five benchmarks (Figure~\ref{fig:overview}; ${\sim}$2.3M binary records), we find a variance structure that reframes what benchmarks actually measure. The dominant source is not the model, the prompt, or the seed---it is the model$\times$item interaction (30--36\%), revealing that models differ primarily in \emph{which items} they find difficult rather than in overall ability. This structure is format-conditional: item-related variance concentrates in multiple-choice benchmarks but disperses in free-form tasks. The practical consequence is immediate: because multi-facet designs average out interaction noise, they achieve reliable measurement with an order of magnitude fewer items than single-condition protocols (Figure~\ref{fig:teaser}).

Our contributions:
\begin{enumerate}
    \item \textbf{Design-dependent reproducibility budgets.} Multi-facet evaluation designs reduce item requirements by up to 10$\times$: a 3-prompt $\times$ 3-seed protocol achieves $G \geq 0.80$ with 178 MMLU items versus 552 under standard conditions. Cross-benchmark budgets range from 93--128 items (multi-facet) versus 1,078--2,768 (single-condition). These budgets complement IRT item selection (Table~\ref{tab:dstudy}).

    \item \textbf{Dominant model$\times$item interaction.} The model$\times$item interaction (30--36\%) exceeds the model main effect by an order of magnitude across all five benchmarks, revealing item-dependent rather than global ability differences. This finding is stable under leave-one-model-out analysis and confirmed by inter-model agreement (Cohen's $\kappa = 0.33$--$0.46$).

    \item \textbf{Format-associated variance patterns (exploratory).} Item-related variance concentrates in multiple-choice benchmarks (68--72\%) but disperses in free-form benchmarks (31--41\%), driven by item difficulty heterogeneity ($\rho = 0.82$). This observation ($n = 5$, $p = 0.10$) warrants broader confirmation.
\end{enumerate}
We release an open-source G-theory pipeline and D-study budget calculator for evaluation design optimization.

\begin{figure}[t]
\centering
\includegraphics[width=\columnwidth]{fig0_teaser_dstudy}
\caption{D-study projections for MMLU (single-model baseline). A recommended 3-prompt $\times$ 3-seed design reduces the budget from 552 to 178 items for $G \geq 0.80$; the full multi-facet design reduces further to 55 items.}
\label{fig:teaser}
\end{figure}
